\clearpage
\VerdierPhysicalPageAnchor{0156}%
\begingroup\small
\setlength{\parskip}{0pt}\setlength{\emergencystretch}{1.5em}

\VerdierSemanticAnchor{ch3-3-2-3-hom-map}
Soient $(X,Y,f)$ et $(X',Y',f')$ deux objets de $\mathrm A(\mathcal D)$.
Pour démontrer la proposition (3.2.3), il suffit de montrer que
l'homomorphisme défini par le foncteur $\operatorname{Im}$
\[
\operatorname{Hom}_{\mathrm A(\mathcal D)}((X,Y,f),(X',Y',f'))
\longrightarrow
\operatorname{Hom}_{\mathcal D^\wedge}(\operatorname{Im}(f),\operatorname{Im}(f'))
\tag{3.2.3.2}
\]
est un isomorphisme.

\VerdierSemanticAnchor{ch3-3-2-3-injectivity}
Montrons d'abord l'injectivité. Soit donc
$(g;g')\colon(X,Y,f)\to(X',Y',f')$ un morphisme de
$\mathcal{F}\ell_1(\mathcal D)$ tel que
$\operatorname{Im}((g;g'))=0$. Le morphisme composé :
\[
\operatorname{Im}((f';\operatorname{id}_{Y'}))\circ
\operatorname{Im}((g;g'))\circ
\operatorname{Im}((\operatorname{id}_X;f))=h(g'f)=0.
\]
Comme $h$ est pleinement fidèle, $g'f=0$ ; donc $(g;g')$ est négligeable.
La surjectivité de l'homomorphisme (3.2.3.2) résultera du lemme suivant.

\VerdierSemanticAnchor{ch3-3-2-4}
\medskip\noindent\textbf{Lemme 3.2.4.}\quad
\textit{Soit
$\alpha\colon\operatorname{Im}(f)\to\operatorname{Im}(f')$ un morphisme de
$\mathcal D^\wedge$. Il existe deux morphismes de $\mathcal D$,
$g\colon X\to X'$ et $g'\colon Y\to Y'$, tels que les diagrammes suivants
soient commutatifs :}

\VerdierSemanticAnchor{ch3-3-2-4-1}
\[
\begin{tikzcd}[column sep=6em,row sep=2.3em]
h(X)\arrow[r,"\operatorname{Im}((\operatorname{id}_X;f))"]\arrow[d,"h(g)"']&
\operatorname{Im}(f)\arrow[d,"\alpha"]\\
h(X')\arrow[r,"\operatorname{Im}((\operatorname{id}_{X'};f'))"']&
\operatorname{Im}(f')
\end{tikzcd}
\tag{3.2.4.1}
\]

\VerdierSemanticAnchor{ch3-3-2-4-2}
\[
\begin{tikzcd}[column sep=5em,row sep=2.3em]
h(X)\arrow[r,"h(f)"]\arrow[d,"h(g)"']&h(Y)\arrow[d,"h(g')"]\\
h(X')\arrow[r,"h(f')"']&h(Y')
\end{tikzcd}
\tag{3.2.4.2}
\]

\VerdierSemanticAnchor{ch3-3-2-4-consequence}
En effet, le lemme implique que
$\alpha=\operatorname{Im}((g;g'))$, par (3.2.3.1) et parce que
$\operatorname{Im}((\operatorname{id}_X;f))$ est un épimorphisme de
$\mathcal D^\wedge$.
\endgroup
